\subsection{The fifth-weight estimate at length two}\label{loc:S17}
Put $P=qf_2$ and $r=Kg$. Since the actual nonnegative factors have
$q+f_2=2+r$, $P\le(1+r/2)^2$.
For $0<d\le1/4$, the intermediate Hermite bound proved in
Section~\ref{loc:S13} is
\[
 \E[p(1+r/2)^8]\le1-d/360.
\]
This is the bound on the auxiliary envelope itself, preceding the final
bound on $P$. Under $K\mapsto-K$, $r=K^2\one$ is unchanged and $p$ becomes
$q$. Thus $\E[qP^4]\le1-d/360$ in this range.

For $1/4\le d\le1/3$, let $c=pq+2r\in[0,2pq]$.
The function $q^5(p^2+c)^4/16$ is convex in $c$; its endpoint chord gives
\[
 qP^4\le B(g)+L(g)c,\quad
 B=p^8q^5/16,\quad L=p^3q^4[(p+2q)^4-p^4]/32.
\]
The intercept maximum occurs at $(p,q)=(16/13,10/13)$, giving $B<9/100$.
For the slope discard the negative fourth power and put
$G(q)=q^4(2-q)^3(2+q)^4/32$.
The logarithm is strictly concave and has derivative
$(16+2q-11q^2)/[q(2-q)(2+q)]$.
Its unique positive zero lies in $(13/10,131/100)$, where the numerator
has endpoint values $1/100,-2571/10000$.
The tangent at $13/10$, whose logarithmic derivative is $10/3003<1/300$,
gives
\[
 G(q_*)\le G(13/10)\frac{30000}{29999}
 <\frac{3631}{1000}\frac{30000}{29999}<91/25.
\]
Here $G(13/10)=11617783760583/3200000000000$.
Since $\E c=1-3d$, the full large interval satisfies
$\E[qP^4]\le9/100+(91/25)(1-3d)\le1$.
We have proved $\E[q^5f_2^4]\le1$ with the actual outward path direction.
The length-one case is $q^5p^4\le q$; multiplication by $E_0$ and full
reversal supply the complete decorated and reversed statements.
